%=========================================================================================
% PACKAGES — ORGANIZED BY CATEGORY
%=========================================================================================

%-----------------------------------------------------------------------------------------
% MATH & SYMBOLS
%-----------------------------------------------------------------------------------------

\usepackage{amsmath}       % Core AMS math features
\usepackage{amsfonts}      % Math fonts
\usepackage{amssymb}       % Extended math symbols
\usepackage{amsthm}        % Theorem and proof environments
\usepackage{mathrsfs}      % Script-style math fonts
\usepackage{yhmath}        % Wide accents
\usepackage{empheq}        % Highlight/emphasise equations
\usepackage{bm}            % Bold math symbols
\usepackage{mathtools}     % Enhancements to amsmath
\usepackage{nccmath}       % Compact math environments 
\usepackage{subcaption}
\usepackage{multicol}
%-----------------------------------------------------------------------------------------
% DOCUMENT LAYOUT
%-----------------------------------------------------------------------------------------

\usepackage{setspace}      % Line spacing control
\usepackage{geometry}      % Page size and margin settings
\geometry{
    a4paper,               % Paper size
    centering,
    left=0.75in,           % Margins
    right=0.75in,
    top=0in+0.75in,
    bottom=0.75in
}

%-----------------------------------------------------------------------------------------
% REFERENCES & LINKS
%-----------------------------------------------------------------------------------------

\usepackage{hyperref}      % Clickable links and cross-references

\hypersetup{
    colorlinks,
    linkcolor={blue!50!black},
    citecolor={blue!50!black},
    urlcolor={blue!80!black},
    pdfpagelayout=TwoPageRight         
}
\usepackage[backend=biber]{biblatex} % Bibliography
% (See https://tug.ctan.org/info/biblatex-cheatsheet/biblatex-cheatsheet.pdf)
\usepackage[nottoc]{tocbibind}   % Add LoF/LoT to TOC
\usepackage{appendix}      % Appendix formatting
\addbibresource{ref.bib}

%-----------------------------------------------------------------------------------------
% TABLES & FIGURES
%-----------------------------------------------------------------------------------------

\usepackage{caption}       % Custom captions
\usepackage{booktabs}      % Professional tables
\usepackage{graphicx}      % Include graphics
\graphicspath{{./Images/}} % Figures directory
\usepackage{multicol}      % Multiple columns
\usepackage{longtable}     % Multi-page tables

%-----------------------------------------------------------------------------------------
% HEADERS / FOOTERS
%-----------------------------------------------------------------------------------------

\usepackage{fancyhdr}      % Custom headers and footers
  \pagestyle{fancy}        % Enable for standard header and footer

%-----------------------------------------------------------------------------------------
% DATE FORMATTING
%-----------------------------------------------------------------------------------------

\usepackage[UKenglish,cleanlook]{isodate}

%-----------------------------------------------------------------------------------------
% COLORS
%-----------------------------------------------------------------------------------------

\usepackage{xcolor}        % Colors
\usepackage{enumitem}
%-----------------------------------------------------------------------------------------
% DRAWING & GRAPHICS
%-----------------------------------------------------------------------------------------

\usepackage{tikz}          % Drawing diagrams
\newcommand\fillin[1][3cm]{\makebox[#1]{\dotfill}} % Dotted Line

%-----------------------------------------------------------------------------------------
% CHEMICAL NOTATION
%-----------------------------------------------------------------------------------------

% \usepackage[version=4]{mhchem} % Chemical formulae and equations
% \usepackage{chemfig}     % Molecule structural formulae
% \usepackage{chemformula} % Chemical notations with more customisation

%-----------------------------------------------------------------------------------------
% OPTIONAL PACKAGES (UNCOMMENT TO ENABLE)
%-----------------------------------------------------------------------------------------

% \usepackage{xwatermark}  % Watermark
% \usepackage{tikz-cd}     % Commutative diagrams
% \usepackage{pgfplots}    % 2D/3D plots
% \usepackage{footmisc}    % Advanced footnotes
% \usepackage{nomencl}     % List of symbols
% \usepackage{makeidx}     % Index
% \usepackage{enumitem}    % Custom lists
% \usepackage{aliascnt}    % Multiple theorem counters
% \usepackage{mdframed}    % Colored/shaded boxes
% \usepackage{psfrag}      % Replace text in EPS
% \usepackage{fontawesome}       % Icon pack (X, FB, ORCID, etc.)
% \usepackage{siunitx}     % Typeset units
% \usepackage{units}       % Typeset units (Legacy)
% \usepackage[english]{babel}    % Language settings
% \usepackage{verbatim}    % Display long code blocks
% \usepackage[l3]{csvsimple}     % CSV table import

%=========================================================================================
% CUSTOM SETTINGS & ENVIRONMENTS
%=========================================================================================

% Chapter and Section Headings in Headers
\renewcommand{\chaptermark}[1]{\markboth{#1}{}}
\renewcommand{\sectionmark}[1]{\markright{\thesection\ #1}}

% Spacing
\setlength{\parskip}{1em plus 0.25em minus 0.25em}
\renewcommand{\baselinestretch}{1.2}

% Theorem Environments
\theoremstyle{plain}
\newtheorem{theorem}{Theorem}[section]
\newtheorem{lemma}[theorem]{Lemma}
\newtheorem{corollary}[theorem]{Corollary}
\newtheorem{proposition}[theorem]{Proposition}
\theoremstyle{definition}
\newtheorem{definition}[theorem]{Definition}
\newtheorem{example}[theorem]{Example}
\newtheorem{notation}[theorem]{Notation}
\theoremstyle{remark}
\newtheorem{remark}[theorem]{Remark}

% Autoref Naming
\providecommand*{\definitionautorefname}{Definition}
\providecommand*{\propositionautorefname}{Proposition}
\providecommand*{\theoremautorefname}{Theorem}
\providecommand*{\corollaryautorefname}{Corollary}
\providecommand*{\lemmaautorefname}{Lemma}

%-----------------------------------------------------------------------------------------
% OTHER DOCUMENT METADATA
%-----------------------------------------------------------------------------------------

\newcommand{\RegNumber}{[Registration Number]}
\newcommand{\Degree}{}
\newcommand{\Department}{[Department Name]}
\newcommand{\ProjectSupervisor}{[Project Supervisor]}
\newcommand{\ThesisCommitteeMemberA}{[TC Member 1]}
\newcommand{\TCMAFacultyPosition}{[Faculty Position]}
\newcommand{\ThesisCommitteeMemberB}{[TC Member 2]}
\newcommand{\TCMBFacultyPosition}{[Faculty Position]}
\newcommand{\ThesisCommitteeMemberC}{[TC Member 3]}
\newcommand{\TCMCFacultyPosition}{[Faculty Position]}

% Hyperref PDF metadata
\newcommand{\ThesisAuthor}{John Doe}
\newcommand{\ThesisTitle}{My Thesis Title}
\newcommand{\ThesisSubject}{\Degree\ Thesis, IISER Thiruvananthapuram}
\newcommand{\ThesisKeywords}{IISER Thiruvananthapuram}

%----------------- Input Commands -----------------%

\newcommand{\SetThesisTitle}[1]{\renewcommand{\ThesisTitle}{#1}}
\newcommand{\SetThesisAuthor}[1]{\renewcommand{\ThesisAuthor}{#1}}
\newcommand{\SetThesisSubject}[1]{\renewcommand{\ThesisSubject}{#1}}
\newcommand{\SetRegNumber}[1]{\renewcommand{\RegNumber}{#1}}
\def\msc{MSc}
\def\phd{PhD}
\def\minor{Minor}
\newcommand{\SetDegree}[1]{%
    \let\Degree=#1%
    \def\DegreeDisplay{}%
    \edef\tempDegree{#1}%
    \ifx\tempDegree\msc
        \def\DegreeDisplay{Master of Science}%
    \else
        \ifx\tempDegree\phd
            \def\DegreeDisplay{Doctor of Philosophy}%
        \else
            \ifx\tempDegree\minor
                \def\DegreeDisplay{Minor Degree}%
            \else
                \def\DegreeDisplay{#1}%
            \fi
        \fi
    \fi
}
\newcommand{\SetDepartment}[1]{\renewcommand{\Department}{#1}}
\newcommand{\SetProjectSupervisor}[1]{\renewcommand{\ProjectSupervisor}{#1}}
\newcommand{\SetThesisCommitteeMemberA}[1]{\renewcommand{\ThesisCommitteeMemberA}{#1}}
\newcommand{\SetTCMAFacultyPosition}[1]{\renewcommand{\TCMAFacultyPosition}{#1}}
\newcommand{\SetThesisCommitteeMemberB}[1]{\renewcommand{\ThesisCommitteeMemberB}{#1}}
\newcommand{\SetTCMBFacultyPosition}[1]{\renewcommand{\TCMBFacultyPosition}{#1}}
\newcommand{\SetThesisCommitteeMemberC}[1]{\renewcommand{\ThesisCommitteeMemberC}{#1}}
\newcommand{\SetTCMCFacultyPosition}[1]{\renewcommand{\TCMCFacultyPosition}{#1}}
\newcommand{\SetThesisKeywords}[1]{\renewcommand{\ThesisKeywords}{#1}}
\usepackage{graphicx}
\usepackage{amsmath, amssymb}
\usepackage{setspace}
\usepackage{titlesec}
\usepackage{xcolor}
\usepackage{url}
\usepackage{hyperref}
\usepackage{longtable}
\usepackage{float}
\usepackage{amsmath,amsfonts,amssymb}
\usepackage{enumitem}
\usepackage{esvect}
\usepackage{graphicx}
\usepackage{pdflscape}
\usepackage{tikz}
\usetikzlibrary{decorations.pathreplacing,decorations.pathmorphing,arrows.meta,calc}
\newcommand{\wavyarrow}{%
  \mathrel{\tikz[baseline=-0.5ex]{
    \draw[->,decorate,
          decoration={snake, amplitude=0.3mm, segment length=2mm}]
          (0,0) -- (1em,0);
  }}%
}

%=========================================================================================
% ABBREVIATIONS, CONSTANTS, SYMBOLS LISTS FOR GLOSSARY
%=========================================================================================

% Abbreviations
\newcommand\listsymbolname{Abbreviations}
\newcommand\listofsymbols[2]{
  \section*{\listsymbolname}
  \addcontentsline{toc}{section}{\listsymbolname}
  \begin{longtable}[c]{#1}
    #2
  \end{longtable}\par
}

% Constants
\newcommand\listconstants{Physical Constants}
\newcommand\listofconstants[2]{
  \section*{\listconstants}
  \addcontentsline{toc}{section}{\listconstants}
  \begin{longtable}[c]{#1}
    #2
  \end{longtable}\par
}

% Symbols (Nomenclature)
\newcommand\listnomenclature{Symbols}
\newcommand\listofnomenclature[2]{
  \section*{\listnomenclature}
  \addcontentsline{toc}{section}{\listnomenclature}
  \begin{longtable}[c]{#1}
    #2
  \end{longtable}\par
}

\setlength{\parindent}{0pt}

\AtEveryBibitem{
  \clearfield{url}
  \clearfield{urlyear}
  \clearfield{urlmonth}
  \clearfield{doi}
  \clearfield{isbn}
  \clearfield{issn}
  \clearfield{eprint}
  \clearfield{eprinttype}
  \clearfield{eprintclass}
}